%%%%%%%%%%%%%%%%%%%%%%%%%%%%%%%%%%%%%%%%%%%%%%%%%%%%%%%%%%%%%%%%%%%%%%
% Lớp 10 | Đại số tổ hợp | Nhị thức Newton
% Dạng 0D8?3-5: Tính tổng nhờ khai triển nhị thức Newton
% Số câu: 40  (sinh tự động bởi tach_cau.py; từ nay thầy sửa trực tiếp file này)
%%%%%%%%%%%%%%%%%%%%%%%%%%%%%%%%%%%%%%%%%%%%%%%%%%%%%%%%%%%%%%%%%%%%%%

%%%%%%%%%%%%%%%% TRẮC NGHIỆM 4 PHƯƠNG ÁN %%%%%%%%%%%%%%%%

% nguồn: MaHoaNH_TracNghiem_Lop10
\begin{ex}%[0D8N3-5]%[ma:3ad079d0]
	Tính tổng các hệ số trong khai  triển $(1-2x)^4$.
	\choice 
		 {\True $1$}
		 {$-1$}
		 {$81$}
		 {$-81$}
	 \loigiai{
		 Tổng các hệ số trong khai triển $(1-2x)^4$ chính là giá trị của biểu thức $(1-2x)^4$ tại $x=1$.\\
		 Vậy $S=(1-2\cdot 1)^4=1$.
	 }
\end{ex}

% nguồn: MaHoaNH_TracNghiem_Lop10
\begin{ex}%[0D8N3-5]%[ma:0653afb7]
	Biết  khai triển nhị thức $P(x)=(2-3 x)^4=a_{4} x^4+a_{3} x^3+a_{2} x^2+a_{1} x+a_{0}$. Tính
	$S=a_{4}+a_{3}+a_{2}+a_{1}+a_{0}$. 
	\choice 
		 {$9$}
		 {$6$}
		 {$3$}
		{\True $1$} 
	 \loigiai{
		Tổng các hệ số $S = a_4 + a_3 + a_2  + a_1 + a_0$ chính là giá trị của đa thức khi thay $x = 1$, tức là
		 $$
		S = P(1)  = (2 - 3 \cdot 1)^4 = (-1)^4 = 1.
		$$ 
	 }
\end{ex}

% nguồn: MaHoaNH_TracNghiem_Lop10
\begin{ex}%[0D8H3-5]%[ma:cea7b135]
	 Biểu thức $P=\mathrm{C}_n^k+\mathrm{C}_n^{k+1}$ bằng
	\choice 
		{\True $\mathrm{C}_{n+1}^{{k+1}}$} 
		{$\mathrm{C}_{n+1}^{k}$} 
		 {$\mathrm{C}_{n+1}^{k}$}
		{$\mathrm{C}_n^{k}$} 
	\loigiai{ 
		Ta có  $C_n^k+C_n^{k+1}=C_{n+1}^{k+1}$.
	 }
\end{ex}

% nguồn: MaHoaNH_TracNghiem_Lop10
\begin{ex}%[0D8H3-5]%[ma:3830a234]
	 Tổng $T=\mathrm{C}_n^0-\mathrm{C}_n^1+\mathrm{C}_n^2+\cdots+(-1)^k \mathrm{C}_n^k+\cdots+(-1)^n \mathrm{C}_n^n$ bằng
	\choice 
		 {$2^{{n+1}}$}
		 {$2^{{n-1}}$}
		 {$2^{n}$}
		{\True  $0$}
	\loigiai{ 
		Theo khai triển  nhị thức Niuton $(a+b)^n=\sum\limits_{k=0}^n C_n^k a^{n-k} b^k$. $\qquad(*)$\\
		Với $a=1; b=-1$, ta có $(*) \Rightarrow 0=C_n^0-C^1+\cdots+(-1)^k  C_n^k+\cdots+(-1)^n C_n^n$.
	} 
\end{ex}

% nguồn: MaHoaNH_TracNghiem_Lop10
\begin{ex}%[0D8H3-5]%[ma:24dc9a00]
	Tổng  $T=\mathrm{C}_{2n}^0+\mathrm{C}_{2n}^2+\mathrm{C}_{2n}^4+\cdots+\mathrm{C}_{2n}^{2k}+\cdots+\mathrm{C}_{2n}^{2n}$ bằng
	\choice 
		 {$2^{{n-1}}$}
		 {$2^{{2n-1}}$}
		{\True $2^{{2n}-1}$} 
		 {$2^{{2n}}$}
	\loigiai{ 
		Ta có  $\mathrm{C}_n^0+\mathrm{C}_n^2+\mathrm{C}_n^4+\cdots=2^{n-1}$.\\
		Áp dụng hệ thức  trên, ta có $T=\mathrm{C}_{2n}^0+\mathrm{C}_{2n}^2+\mathrm{C}_{2n}^4+\cdots+\mathrm{C}_{2n}^{2k}+\cdots+\mathrm{C}_{2n}^{2n}=2^{2n-1}$.
	 }
\end{ex}

% nguồn: MaHoaNH_TracNghiem_Lop10
\begin{ex}%[0D8H3-5]%[ma:7837a186]
	Tổng $T=\mathrm{C}_n^0+\mathrm{C}_n^1+\mathrm{C}_n^3+\mathrm{C}_n^4+\cdots..+\mathrm{C}_n^n$  bằng
	 \choice
		{$2^{{n+1}}$} 
		 {$2^{{n-1}}$}
		{\True  $2^{n}$}
		{$0$} 
	\loigiai{ 
		Theo khai triển  nhị thức Niuton $(a+b)^n=\sum\limits_{k=0}^n \mathrm{C}_n^k a^{n-k} b^k$. $\qquad(*)$\\
		Với $a=b=1$, ta có $(*) \Rightarrow 2^n=\mathrm{C}_n^0+C_n^1+\cdots+\mathrm{C}_n^{n-1}+\mathrm{C}_n^n$. 
	 }
\end{ex}

% nguồn: MaHoaNH_TracNghiem_Lop10
\begin{ex}%[0D8H3-5]%[ma:9d0e8e2d]
	 Tính giá trị của $H=\mathrm{C}_{15}^0-2\mathrm{C}_{15}^1+2^2\mathrm{C}_{15}^2-\cdots+2^{14} \mathrm{C}_{15}^{14}-2^{15} \mathrm{C}_{15}^{15}$.
	 \choice
		 {$-3^{{15}}$}
		 {$3^{{15}}$}
		{$1$} 
		{\True  $-1$}
	\loigiai{ 
		$(1+x)^{15}=\mathrm{C}_{15}^0+\mathrm{C}_{15}^1x+\mathrm{C}_{15}^2x^2+\cdots+\mathrm{C}_{15}^{14} x^{14}+\mathrm{C}_{15}^{15} x^{15}$.\\ 
		Chọn $x=-2$,  ta được $\mathrm{C}_{15}^0-2\mathrm{C}_{15}^1+2^2\mathrm{C}_{15}^2-\cdots+2^{14} \mathrm{C}_{15}^{14}-2^{15} \mathrm{C}_{15}^{15}=(1-2)^{15}=-1$.
	} 
\end{ex}

% nguồn: MaHoaNH_TracNghiem_Lop10
\begin{ex}%[0D8H3-5]%[ma:afc9ecba]
	Xét khai triểm $\left(1+2x+x^2\right)^{20}=a_0+a_1 x+\cdots+a_{40}  x^{40}$. Tổng $S=a_0+a_1+\cdots+a_{40}$ là
	 \choice
		{$4^{{40}}$} 
		{$2^{{20}}$} 
		{\True $2^{{40}}$} 
		 {$4^{{10}}$}
	\loigiai{ 
		Xét  khai triển $\left(1+2x+x^2\right)^{20}=(1+x)^{40}=C_{40}^0+C_{40}^1x+C_{40}^2x^2+\cdots+C_{40}^{40} x^{40}$.\\
		Chọn $x=1$ ta được $S=a_0+a_1+\cdots+a_{40}=2^{40}$. 
	} 
\end{ex}

% nguồn: MaHoaNH_TracNghiem_Lop10
\begin{ex}%[0D8H3-5]%[ma:cd814269]
	Tính tổng các hệ số trong khai triển nhị thức Newtơn  của $(1-2x)^4$.
	 \choice
		{\True  $1$}
		 {$-1$}
		 {$81$}
		 {$-81$}
	\loigiai{ 
		Khai triển  $(2x-3)^4$ có dạng $(2x-3)^4=a_0x^4+a_1x^3+a_2x^2+a_3x+a_4$ \hfill $(1)$\\
		Từ $(1)$, cho  $x=1$ ta được $a_0+a_1+a_2+a_3+a_4=(2\cdot 1-3)^4=1$.\\
		Vậy $S=1$. 
	} 
\end{ex}

% nguồn: MaHoaNH_TracNghiem_Lop10
\begin{ex}%[0D8H3-5]%[ma:2a00f89c]
	Cho $T=\mathrm{C}_{2\,022}^1+\mathrm{C}_{2\,022}^3+\mathrm{C}_{2\,022}^5+\cdots+\mathrm{C}_{2\,022}^{2\,021}$. Tính biểu thức $T=2^n$  thì $n$ bằng
	\choice 
		{$2\,023$} 
		 {$2\,022$}
		{\True $2\,021$} 
		{$2\,020$} 
	 \loigiai{
		Ta  có $\mathrm{C}_n^1+\mathrm{C}_n^3+\mathrm{C}_n^5+\cdots+\mathrm{C}_n^n=2^{n-1}$.\\
		Áp dụng  $T=\mathrm{C}_{2\,022}^1+\mathrm{C}_{2\,022}^3+\mathrm{C}_{2\,022}^5+\cdots+\mathrm{C}_{2\,022}^{2\,021}=2^{2\,021}$.\\
		 Do đó $n=2\,021$.
	} 
\end{ex}

% nguồn: MaHoaNH_TracNghiem_Lop10
\begin{ex}%[0D8H3-5]%[ma:6a3a8faf]
	Cho $n$ là số nguyên dương thỏa mãn $\mathrm{C}_n^1+\mathrm{C}_n^2+\cdots+\mathrm{C}_n^n=4\,095$. Giá trị  của n bằng
	\choice 
		 {$14$}
		 {$16$}
		 {$13$}
		 {\True $12$}
	 \loigiai{
		Ta có $\mathrm{C}_n^1+\mathrm{C}_n^2+\cdots+\mathrm{C}_n^n=4\,095\Rightarrow \mathrm{C}_n^0+\mathrm{C}_n^1+\mathrm{C}_n^2+\cdots+\mathrm{C}_n^n=4\,096$.\\ 
		Mà $\mathrm{C}_n^0+\mathrm{C}_n^1+\mathrm{C}_n^2+\cdots+\mathrm{C}_n^n=2^n$  nên suy ra
		$2^n=4\,096\Leftrightarrow n=12$. 
	 }
\end{ex}

% nguồn: MaHoaNH_TracNghiem_Lop10
\begin{ex}%[0D8H3-5]%[ma:ac03341a]
	Tính giá  trị của $K=3^{20} \mathrm{C}_{20}^0-3^{19} \cdot 4\cdot \mathrm{C}_{20}^1+3^{18} \cdot 4^2\cdot \mathrm{C}_{20}^2-\cdots-3\cdot 4^{19} \cdot \mathrm{C}_{20}^{19}+4^{20} \cdot \mathrm{C}_{20}^{20}$.
	 \choice
		 {$7^{{20}}$}
		{$-7^{{20}}$} 
		{$-1$} 
		 {\True $1$}
	\loigiai{ 
		$(3+x)^{20}=3^{20} \mathrm{C}_{20}^\circ+3^{19} \mathrm{C}_{20}^1x+3^{18}  \mathrm{C}_{20}^2x^2+\cdots+3\mathrm{C}_{20}^{19} x^{19}+\mathrm{C}_{20}^{20} x^{20}$.\\
		Chọn $x=-4$ ta được 
		\[3^{20}  \mathrm{C}_{20}^0-3^{19} \cdot 4\cdot \mathrm{C}_{20}^1+3^{18} \cdot 4^2\cdot \mathrm{C}_{20}^2-\cdots-3.4^{19} \cdot \mathrm{C}_{20}^{19}+4^{20} \cdot \mathrm{C}_{20}^{20}=(3-4)^{20}=1.\]
	} 
\end{ex}

% nguồn: MaHoaNH_TracNghiem_Lop10
\begin{ex}%[0D8H3-5]%[ma:b2d42ada]
	Tính tổng $\mathrm{C}_n^0-\mathrm{C}_n^1+\mathrm{C}_n^2+\cdots+(-1)^n \mathrm{C}_n^n$. ta  được kết quả là:
	\choice 
		{\True  $0$}
		 {$2^{n}$}
		{$2^{{n-1}}$} 
		 {$2^{{n+1}}$}
	 \loigiai{
		Xét khai triển: $(a+b)^n=C_n^0a^n+C_n^1a^{n-1}  b+C_n^2a^{n-2} b^2+\cdots+C_n^n b^n$.\\
		Chọn $\left\{\begin{array}{l}a=1\\ b=-1
	\end{array}\right.$ ta được $(1-1)^n=C_n^0\cdot 1^n+C_n^1\cdot 1^{n-1} \cdot(-1)+C_n^2\cdot 1^{n-2} \cdot(-1)^2+\cdots+C_n^n \cdot(-1)^n$
	 $\Leftrightarrow 0=C_n^0-C_n^1+C_n^2+\cdots+(-1)^n C_n^n$.
	 }
\end{ex}

% nguồn: MaHoaNH_TracNghiem_Lop10
\begin{ex}%[0D8H3-5]%[ma:ac4dab5a]
	 Với $n \geq 4$, tổng $T=\mathrm{C}_n^0+\mathrm{C}_n^2+\mathrm{C}_n^4+\cdots$ bằng
	\choice 
		{$2^{{2n-1}}$} 
		{\True  $2^{{n-1}}$}
		 {$2^{n}$}
		{$2^n-1$} 
	\loigiai{ 
		 Theo khai triển nhị thức Niuton $(a+b)^n=\sum\limits_{k=0}^n \mathrm{C}_n^k a^{n-k} b^k$. $\qquad(*)$\\
		Với $a=b=1$, ta có $(*)  \Rightarrow 2^n=\mathrm{C}_n^0+C_n^1+\cdots+\mathrm{C}_n^{n-1}+\mathrm{C}_n^n$. $\qquad(1)$\\
		Với $a=1$, $b=-1$, ta có  $(*) \Rightarrow 0=\mathrm{C}_n^0-C_n^1+\cdots+(-1)^n\mathrm{C}_n^n$. $\qquad(2)$\\
		Lấy $(1)+(2)$ suy  ra $2^n=2T\Rightarrow T=2^{n-1}$.
	 }
\end{ex}

% nguồn: MaHoaNH_TracNghiem_Lop10
\begin{ex}%[0D8H3-5]%[ma:69b4f9f3]
	Thu  gọn biểu thức $A=\left(2+\sqrt{3}\right)^5-\left(2-\sqrt{3}\right)^5$ ta được $A=a+b \sqrt{3}$ với $a$, $b$ là các số nguyên. Tính tổng $a+b$.
	\choice 
		 {$-209$}
		 {$209$}
		{$-418$} 
		{\True $418$} 
	\loigiai{ 
		 Ta có\\
		\indent $\left(2+\sqrt{3}\right)^5=\mathrm{C}_{5}^0 2^5+\mathrm{C}_{5}^1 2^4 \sqrt{3}+\mathrm{C}_{5}^2 2^3\left(\sqrt{3}\right)^2+\mathrm{C}_{5}^3 2^2\left(\sqrt{3}\right)^3+\mathrm{C}_{5}^4  2\left(\sqrt{3}\right)^4+\mathrm{C}_{5}^5\left(\sqrt{3}\right)^5$.\\
		\indent $\left(2-\sqrt{3}\right)^5=\mathrm{C}_{5}^0 2^5-\mathrm{C}_{5}^1 2^4 \sqrt{3}+\mathrm{C}_{5}^2 2^3\left(\sqrt{3}\right)^2-\mathrm{C}_{5}^3 2^2\left(\sqrt{3}\right)^3+\mathrm{C}_{5}^4  2\left(\sqrt{3}\right)^4-\mathrm{C}_{5}^5\left(\sqrt{3}\right)^5$.\\
		Do đó $A=2\left[\mathrm{C}_{5}^1  2^4 \sqrt{3}+\mathrm{C}_{5}^3 2^2\left(\sqrt{3}\right)^3+\mathrm{C}_{5}^5\left(\sqrt{3}\right)^5\right]=418 \sqrt{3}$.\\
		 Vậy $a+b=418$.
	} 
\end{ex}

% nguồn: MaHoaNH_TracNghiem_Lop10
\begin{ex}%[0D8H3-5]%[ma:5b79c59a]
	Cho khai triển  $\left(7-8x\right)^5=a_0+a_1x+a_2x^2+a_3x^3+a_4x^4+a_5x^5$. Tính giá trị của biểu thức $S=a_0+a_1+a_2+a_3+a_4+a_5$.
	 \choice
		 {\True $S=-1$}
		{$S=0$} 
		 {$S=2$}
		{$S=1$} 
	\loigiai{ 
	Ta có 
	\begin{eqnarray*}
		S &=&  a_0+a_1+a_2+a_3+a_4+a_5\\
		&=& a_0+a_1 \cdot 1+a_2 \cdot 1^2+a_3 \cdot 1^3+a_4 \cdot 1^4+a_5 \cdot  1^5 = (7-8\cdot 1)^5= (-1)^5=-1.
	\end{eqnarray*}
	} 
\end{ex}

% nguồn: MaHoaNH_TracNghiem_Lop10
\begin{ex}%[0D8H3-5]%[ma:2806896d]
	Tính tổng $\mathrm{C}_{2n}^0+\mathrm{C}_{2n}^2+\mathrm{C}_{2n}^4+\cdots+\mathrm{C}_{2n}^{2n}$ ta được  kết quả là
	\choice 
		{$2^{{n-1}}$} 
		{$2^{n}$} 
		{\True  $2^{{2n-1}}$}
		 {$2^{{2n+1}}$}
	 \loigiai{
		Xét  khai triển $(a+b)^{2n}=C_{2n}^0a^{2n}+C_{2n}^1a^{2n-1} b+C_{2n}^2a^{2n-2} b^2+\cdots+C_{2n}^{2n} b^{2n}$.\\
		Chọn $\left\{\begin{array}{l}a=1\\ b=1
	\end{array}\right.$ ta được $2^{2n}=C_{2n}^0+C_{2n}^1+C_{2n}^2+\cdots+C_{2n}^{2n}$. $\qquad(1)$\\
	Chọn $\left\{\begin{array}{l}a=1\\ b=-1
	\end{array}\right.$ ta được $0=C_{2n}^0-C_{2n}^1+C_{2n}^2-C_{2n}^3+C_{2n}^4+\cdots-C_{2n}^{2n-1}+C_{2n}^{2n}$. $\qquad(2)$\\
	Từ $(1)$ và $(2)$  suy ra $C_{2n}^0+C_{2n}^2+C_{2n}^4+\cdots+C_{2n}^{2n}=2^{2n-1}$.
	} 
\end{ex}

% nguồn: MaHoaNH_TracNghiem_Lop10
\begin{ex}%[0D8H3-5]%[ma:408a8b8a]
	Tính tổng $\mathrm{C}_n^0+\mathrm{C}_n^1+\mathrm{C}_n^2+\cdots+\mathrm{C}_n^n$.  ta được kết quả là
	 \choice
		 {$3^{n}$}
		 {$2^{n}$}
		 {$n !$}
		{$2^{{n+1}}$} 
	\loigiai{ 
		Xét khai  triển $(a+b)^n=\mathrm{C}_n^0a^n+\mathrm{C}_n^1a^{n-1} b+\mathrm{C}_n^2a^{n-2} b^2+\cdots+\mathrm{C}_n^n b^n$.\\
		Chọn $\left\{\begin{array}{l}a=1\\ b=1
	\end{array}\right.$ ta được $(1+1)^n=\mathrm{C}_n^0\cdot 1^n+\mathrm{C}_n^1\cdot 1^{n-1} \cdot 1+\mathrm{C}_n^2\cdot 1^{n-2} \cdot 1^2+\cdots+\mathrm{C}_n^n \cdot 1^n$
	$$ 
	\Leftrightarrow 2^n=\mathrm{C}_n^0+\mathrm{C}_n^1+\mathrm{C}_n^2+\cdots+\mathrm{C}_n^n. 
	 $$
	 }
\end{ex}

% nguồn: MaHoaNH_TracNghiem_Lop10
\begin{ex}%[0D8H3-5]%[ma:ee6d2ce0]
	Tính tổng $\mathrm{C}_n^1+2\dfrac{\mathrm{C}_n^2}{\mathrm{C}_n^1}+3\dfrac{\mathrm{C}_n^3}{\mathrm{C}_n^2}+\cdots.+n \dfrac{\mathrm{C}_n^n}{\mathrm{C}_n^{n-1}}$ ta được kết quả là 
	\choice 
		 {$3^{n}$}
		{$2^{n}$} 
		 {$\dfrac{n(n-1)}{2}$}
		{\True  $\dfrac{n(n+1)}{2}$}
	 \loigiai{
		 Ta có $\dfrac{\mathrm{C}_n^k}{\mathrm{C}_n^{k-1}}=\dfrac{n-k+1}{k}$.\\
		 Suy ra
		 \allowdisplaybreaks
		\begin{align*}
			T&=\mathrm{C}_n^1+2\dfrac{\mathrm{C}_n^2}{\mathrm{C}_n^1}+3\dfrac{\mathrm{C}_n^3}{\mathrm{C}_n^2}+\cdots.+n  \dfrac{\mathrm{C}_n^n}{\mathrm{C}_n^{n-1}}
			\\ 
			&=n+2\cdot \dfrac{n-1}{2}+3\dfrac{n-2}{3}+\cdots+n \cdot \dfrac{1}{n} 
			 \\
			&=n+(n-1)+(n-2)+\cdots+2+1=\dfrac{n(n+1)}{2}. 
		\end{align*}
	} 
\end{ex}

% nguồn: MaHoaNH_TracNghiem_Lop10
\begin{ex}%[0D8H3-5]%[ma:30da92c9]
	Tính tổng $n \cdot 2^{n-1} \cdot \mathrm{C}_n^0+(n-1) \cdot 2^{n-2}  \cdot 3\cdot \mathrm{C}_n^1+(n-2) \cdot 2^{n-3} \cdot 3^2\cdot \mathrm{C}_n^2+\cdots+3^{n-1} \cdot \mathrm{C}_n^{n-1}$ ta được kết quả là
	 \choice
		{$5^{n}$} 
		{$n{,}5^{n}$} 
		 {\True $n{,}5^{{n-1}}$}
		 {$5^{{n-1}}$}
	\loigiai{ 
		Ta  có
		\allowdisplaybreaks 
		\begin{align*}
			T&=n \cdot 2^{n-1} \cdot  C_n^0+(n-1) \cdot 2^{n-2} \cdot 3\cdot C_n^1+(n-2) \cdot 2^{n-3} \cdot 3^2\cdot C_n^2+\cdots+3^{n-1} \cdot C_n^{n-1}
			 \\
			&=\sum_{k=0}^{n-1}(n-k) \cdot 2^{n-k-1}  \cdot 3^k \cdot C_n^k
			\\ 
			 &=\sum_{k=0}^{n-1} n \cdot 2^{n-k-1} \cdot 3^k \cdot C_{n-1}^{n-k-1}=n \cdot(2+3)^{n-1}
			\\ 
			 &=n \cdot 5^{n-1}.
		\end{align*}
	 }
\end{ex}

% nguồn: MaHoaNH_TracNghiem_Lop10
\begin{ex}%[0D8H3-5]%[ma:00988c73]
	Cho $n$ là số nguyên dương thỏa mãn  $\mathrm{C}_n^7+\mathrm{C}_n^8=\mathrm{C}_{n+1}^9$. Giá trị của số n bằng
	\choice 
		 {\True $16$}
		{$24$} 
		 {$18$}
		 {$17$}
	\loigiai{ 
		Điều kiện  $n \geq 8; n \in \mathbb{N}$.\\
		 Áp dụng $C_n^k+C_n^{k+1}=C_{n+1}^{k+1}$.\\
		 Ta có
		 \allowdisplaybreaks
		\begin{align*}
			&C_n^7+C_n^8=C_{n+1}^9 
			\\ 
			\Leftrightarrow& C_{n+1}^8=C_{n+1}^9 
			 \\
			\Leftrightarrow& \dfrac{(n+1)!}{8!(n-7)!}=\dfrac{(n+1)!}{9!(n-8)!} 
			\\ 
			 \Leftrightarrow& \dfrac{1}{n-7}=\dfrac{1}{9}
			 \\
			\Leftrightarrow&  n=16.
		\end{align*}
	} 
\end{ex}

% nguồn: MaHoaNH_TracNghiem_Lop10
\begin{ex}%[0D8H3-5]%[ma:be5c0d4b]
	Cho khai triển $(3x-5)^n=\sum\limits_{i=0}^n a_i x^i$. Biết $\mathrm{C}_n^0+2\mathrm{C}_n^1+4\mathrm{C}_n^2+\cdots+2^n \mathrm{C}_n^n=243$.  Khi đó hãy tính tổng $S=a_0+a_1+a_2+\cdots+a_{n-1}$.
	 \choice
		{\True  $3\,093$}
		{$-3\,157$} 
		{$3\,157$} 
		 {$-3\,093$}
	\loigiai{ 
		Ta có $\mathrm{C}_n^0+2\mathrm{C}_n^1+4\mathrm{C}_n^2+\cdots+2^n \mathrm{C}_n^n=243\Leftrightarrow(1+2)^n=243\Leftrightarrow  3^n=3^5\Leftrightarrow n=5$.\\
		Ta có 
		\allowdisplaybreaks 
		\begin{align*}
			 f(x)&=(3x-5)^5
			\\ 
			&=\mathrm{C}_5^0(3x)^5+\mathrm{C}_5^1(3x)^4(-5)+\mathrm{C}_5^2(3x)^3(-5)^2+\mathrm{C}_5^3(3x)^2(-5)^3+\mathrm{C}_5^4(3x)(-5)^4+\mathrm{C}_5^5(-5)^5. 
		\end{align*}
		Tổng là 
		 \allowdisplaybreaks
		\begin{align*}
			S&=\mathrm{C}_5^03^5+\mathrm{C}_5^13^4(-5)+\mathrm{C}_5^23^3(-5)^2+\mathrm{C}_5^33^2(-5)^3+\mathrm{C}_5^4\cdot 3\cdot(-5)^4 
			 \\
			 &=f(1)-\mathrm{C}_5^5(-5)^5
			\\ 
			 &=(3-5)^5+5^5=3\,093.
		\end{align*}
	 }
\end{ex}

% nguồn: MaHoaNH_TracNghiem_Lop10
\begin{ex}%[0D8H3-5]%[ma:15f857b9]
	Cho $n$ là số nguyên dương thỏa mãn  $\mathrm{C}_{n+4}^{n+1}-\mathrm{C}_{n+3}^n=8(n+2)$.
	 \choice
		{$14$} 
		{\True $13$} 
		 {$16$}
		 {$15$}
	 \loigiai{
		 Điều kiện $n \in \mathbb{N}$.\\
		 Ta có
		 \allowdisplaybreaks
		\begin{align*}
			 &C_{n+4}^{n+1}-C_{n+3}^n=8(n+2)
			\\ \Leftrightarrow&\left(C_{n+3}^n+C_{n+3}^{n+1}\right)-C_{n+3}^n=8(n+2) 
			 \\
			\Leftrightarrow& C_{n+3}^{n+1}=8(n+2) 
			 \\
			\Leftrightarrow&  \dfrac{(n+2)(n+3)}{2!}=8(n+2)
			 \\
			 \Leftrightarrow &n+3=8\cdot 2!
			\\ 
			\Leftrightarrow  &n+3=16
			\\ 
			\Leftrightarrow&  n=13.
		\end{align*}
	 }
\end{ex}


%%%%%%%%%%%%%%%% ĐÚNG/SAI %%%%%%%%%%%%%%%%

% nguồn: MaHoaNH_DungSai_Lop10
\begin{ex}%[0D8H3-5]%[ma:f458f995]
	Trong khai triển nhị thức Newton của  $\left(xy+2\right)^5$.
	 \choiceTF
		{Biểu thức triển khai có $5$ số  hạng}
		{\True Số hạng  thứ $3$ là $40x^3y^3$}
		{Hệ số chứa $x^2y^2$ trong khai triển  là $80$}
		{\True Tổng hệ số của  các số hạng trong khai triển là $243$}
	\loigiai{ 
		\begin{itemchoice}
			\itemch Do $n=5$ nên khai triển có $6$ số hạng. 
			\itemch $\left(xy+2\right)^5  = x^5y^5 + 10x^4y^4 + 40x^3y^3 + 80x^2y^2 + 80xy + 32$.
			 \itemch Số hạng chứa $x^2y^2$ là $80x^2y^2$.
			\itemch Áp dụng  công thức ta có $1 + 10 + 40 + 80 + 80 + 32 = 243$.
		\end{itemchoice}
	 }
\end{ex}

% nguồn: MaHoaNH_DungSai_Lop10
\begin{ex}%[0D8H3-5]%[ma:d889c112]
	Khai  triển $(1-x)^6$.
	\choiceTF 
		 {\True Hệ số của $x^2$ trong khai triển là $\mathrm{C}_6^2$}
		{Hệ số của $x^3$ trong khai triển là $\mathrm{C}_6^3$} 
		{\True Hệ số của $x^5$ trong khai triển  là $-\mathrm{C}_6^5$}
		 {$\mathrm{C}_6^0-\mathrm{C}_6^1+\mathrm{C}_6^2-\mathrm{C}_6^3+\mathrm{C}_6^4-\mathrm{C}_6^5+\mathrm{C}_6^6=1$}
	\loigiai{Ta có 
		\begin{align*}
			(1-x)^6 &= \mathrm{C}_6^0 - \mathrm{C}_6^1x + \mathrm{C}_6^2x^2 -  \mathrm{C}_6^3x^3 + \mathrm{C}_6^4x^4 - \mathrm{C}_6^5x^5 + \mathrm{C}_6^6x^6
			 \\&=1-6x+15x^2-20x^3+15x^4-6x^5+x^6
			\\&=x^6 - 6x^5 + 15x^4  - 20x^3 + 15x^2 - 6x + 1.\qquad (*)
		\end{align*}
		\begin{itemchoice}
			\itemch Hệ số của $x^2$  trong khai triển là $\mathrm{C}_6^2$.
			 \itemch Hệ số của $x^3$ trong khai triển là $-\mathrm{C}_6^3$.
			\itemch Hệ số của  $x^5$ trong khai triển là $-\mathrm{C}_6^5$.
			 \itemch Thay $x=1$ vào $(*)$, ta được
			$(1-1)^6 = \mathrm{C}_6^0 - \mathrm{C}_6^1  + \mathrm{C}_6^2 - \mathrm{C}_6^3 + \mathrm{C}_6^4 - \mathrm{C}_6^5 + \mathrm{C}_6^6 =0$.
		\end{itemchoice}
	 }
\end{ex}

% nguồn: MaHoaNH_DungSai_Lop10
\begin{ex}%[0D8H3-5]%[ma:1f329e73]
	Khai triển  $\left(x+2y\right)^3+\left(2x-y\right)^3$.
	 \choiceTF
		{\True Hệ  số của $x^3$ là $9$}
		{\True Hệ  số của $y^3$ là $7$}
		{Hệ số của $x^2y$  là $6$}
		{Tổng các hệ số của số hạng mà lũy thừa  của $x$ lớn hơn lũy thừa của $y$ bằng $-3$}
	\loigiai{Ta  có
		\begin{align*}
			\left(x+2y\right)^3 + \left(2x-y\right)^3 
			&= \mathrm{C}_3^0 x^3 \left(2y\right)^0 + \mathrm{C}_3^1 x^2 \left(2y\right)^1 + \mathrm{C}_3^2 x^1 \left(2y\right)^2  + \mathrm{C}_3^3 x^0 \left(2y\right)^3
			 \\&+ \mathrm{C}_3^0 \left(2x\right)^3 \left(-y\right)^0 + \mathrm{C}_3^1 \left(2x\right)^2 \left(-y\right)^1 + \mathrm{C}_3^2 \left(2x\right)^1 \left(-y\right)^2 + \mathrm{C}_3^3 \left(2x\right)^0 \left(-y\right)^3
			\\&= x^3 + 6x^2y + 12xy^2 + 8y^3 +  8x^3 - 12x^2y + 6xy^2 - y^3
			\\&= 9x^3 - 6x^2y + 18xy^2 + 7y^3. 
		\end{align*}
		\begin{itemchoice}
			\itemch Hệ số của  $x^3$ là $9$.
			\itemch Hệ số  của $y^3$ là $7$.
			\itemch Hệ số  của $x^2y$ là $-6$.
			\itemch Tổng các hệ số của số hạng mà lũy thừa của $x$ lớn hơn lũy  thừa của $y$ bằng $3$.
		\end{itemchoice}
	 }
\end{ex}

% nguồn: MaHoaNH_DungSai_Lop10
\begin{ex}%[0D8H3-5]%[ma:7c9ed292]
	Cho  $\left(1-\dfrac{1}{2}x\right)^5=a_0+a_1x+a_2x^2+a_3x^3+a_4x^4+a_5x^5$.
	\choiceTF 
		{$a_3=\dfrac{5}{2}$} 
		{\True $a_5=-\dfrac{1}{32}$} 
		{\True Hệ số lớn nhất trong tất cả hệ số  là $\dfrac{5}{2}$}
		{Tổng $a_0+a_1+a_2+a_3+a_4+a_5=\dfrac{1}{16}$} 
	\loigiai{ 
		\begin{align*}
			 \left(1-\dfrac{1}{2}x\right)^5
			&= \mathrm{C}_5^0 + \mathrm{C}_5^1\left(-\dfrac{1}{2}x\right) +  \mathrm{C}_5^2\left(-\dfrac{1}{2}x\right)^2 + \mathrm{C}_5^3\left(-\dfrac{1}{2}x\right)^3 + \mathrm{C}_5^4\left(-\dfrac{1}{2}x\right)^4 + \mathrm{C}_5^5\left(-\dfrac{1}{2}x\right)^5
			 \\&=1-\dfrac{5}{2}x+\dfrac{5}{2}x^2-\dfrac{5}{4}x^3+\dfrac{5}{16}x^4-\dfrac{1}{32}x^5
			\\&=  a_0 + a_1x + a_2x^2 + a_3x^3 + a_4x^4 + a_5x^5.\qquad (*)
		\end{align*}
		Suy ra $a_0=1,  a_1=-\dfrac{5}{2}, a_2=\dfrac{5}{2}, a_3=-\dfrac{5}{4}, a_4=\dfrac{5}{16}, a_5=-\dfrac{1}{32}$.
		\begin{itemchoice}
			\itemch  Ta có $a_3=-\dfrac{5}{4}$.
			\itemch Ta có $a_5=-\dfrac{1}{32}$. 
			\itemch Ta thấy hệ số lớn nhất tìm được là $a_2=\dfrac{5}{2}$. 
			\itemch Thay $x=1$ vào $(*)$, ta được  $\left(1-\dfrac{1}{2}\right)^5 = a_0 + a_1 + a_2 + a_3 + a_4 + a_5$.\\ Vậy $a_0 + a_1 + a_2 + a_3 + a_4 + a_5 = \dfrac{1}{32}$.
		\end{itemchoice}
	} 
\end{ex}

% nguồn: MaHoaNH_DungSai_Lop10
\begin{ex}%[0D8H3-5]%[ma:fd453871]
	Khai triển  $\left(x+2y\right)^{3}+\left(2x-y\right)^{3}$.
	 \choiceTF
		{\True Hệ  số của $x^3$ là $9$}
		{\True Hệ số của $y^3$  là $7$}
		{Hệ số của $x^2y$ là $6$} 
		{Tổng  các hệ số của số hạng mà lũy thừa của $x$ lớn hơn lũy thừa của $y$ bằng-3}
	\loigiai{Ta có 
		\begin{align*}
			 \left(x+2y\right)^3 + \left(2x-y\right)^3 &
			= \mathrm{C}_3^0x^3 + \mathrm{C}_3^1x^2\left(2y\right) + \mathrm{C}_3^2x\left(2y\right)^2 + \mathrm{C}_3^3\left(2y\right)^3 \\&+ \mathrm{C}_3^0\left(2x\right)^3 +  \mathrm{C}_3^1\left(2x\right)^2\left(-y\right) + \mathrm{C}_3^2\left(2x\right) \left(-y\right)^2 + \mathrm{C}_3^3\left(-y\right)^3
			 \\&= x^3 + 6x^2y + 12xy^2 + 8y^3 + 8x^3 - 12x^2y + 6xy^2 - y^3
			\\&= 9x^3 - 6x^2y + 18xy^2 + 7y^3 
		\end{align*}
		\begin{itemchoice}
			\itemch Hệ số của $x^3$ trong khai triển là $9$. 
			\itemch Hệ số của $y^3$  trong khai triển là $7$.
			\itemch Hệ  số của $x^2y$ trong khai triển là -$6$.
			\itemch Có hai số hạng  mà lũy thừa của $x$ lớn hơn lũy thừa của $y$ là $9x^3, -6x^2y$. \\
			Tổng hệ số của chúng là $9 + \left(-6\right)  = 3$.
		\end{itemchoice}
	 }
\end{ex}

% nguồn: MaHoaNH_DungSai_Lop10
\begin{ex}%[0D8V3-5]%[ma:07fd9a29]
	Chọn đẳng  thức đúng đúng trong các đẳng thức sau
	\choiceTF 
		{\True  $S=\mathrm{C}_{5}^{0}+2\mathrm{C}_{5}^{1}+2^{2}\mathrm{C}_{5}^{2}+ \cdots+2^5\mathrm{C}_{5}^{5}=3^5$}
		{\True $S=4^0\mathrm{C}_{8}^{0}+4^1\mathrm{C}_{8}^{1}+4^{2}\mathrm{C}_{8}^{2}+ \cdots  +4^8\mathrm{C}_{8}^{8}=5^8$}
		 {$3^{10}\mathrm{C}_{10}^{0}-3^{9}\mathrm{C}_{10}^{1}+3^{8}\mathrm{C}_{10}^{2}-3^{7}\mathrm{C}_{10}^{3} +\cdots+3^{1}\mathrm{C}_{10}^{9}-\mathrm{C}_{10}^{10}=2048$}
		{$\mathrm{C}_{12}^{1}+\mathrm{C}_{12}^{2} +\cdots+\mathrm{C}_{12}^{11}  +\mathrm{C}_{12}^{12}=4096$}
	\loigiai{ 
		\begin{itemchoice}
			\itemch  Ta có
			 $(a+b)^{5}=\mathrm{C}_{5}^{0}a^{5}+\mathrm{C}_{5}^{1}a^{4}b+\mathrm{C}_{5}^{2}a^{3}b^{2}+\mathrm{C}_{5}^{3}a^{2}b^3 +\mathrm{C}_{5}^{4}ab^{4}+\mathrm{C}_{5}^{5}b^{5}$.\\
			Cho $a=1$, $b=2$,  ta có
			\[3^{5}=\mathrm{C}_{5}^{0}+2\mathrm{C}_{5}^{1}+2^2\mathrm{C}_{5}^{2}+  2^3\mathrm{C}_{5}^{3}+2^4\mathrm{C}_{5}^{4}+2^5\mathrm{C}_{5}^{5}.\]
			 Vậy $S=3^5$.
			 \itemch Ta có
			$(a+b)^{8}=\mathrm{C}_{8}^{0}a^{8}b^0+\mathrm{C}_{8}^{1}a^{7}b^1+\mathrm{C}_{8}^{2}a^{6}b^{2}+ \cdots + \mathrm{C}_{8}^{8}b^8$.\\ 
			Cho $a=1$, $b=4$, ta có 
			\[5^{8}=4^0\mathrm{C}_{8}^{0}+4^1\mathrm{C}_{8}^{1}+4^2\mathrm{C}_{8}^{2}+ \cdots  +4^8\mathrm{C}_{8}^{8}.\]
			 Vậy $S=5^8$.
			\itemch  Ta có
			$(a+b)^{n}=\mathrm{C}_{n}^{0}a^{n}+\mathrm{C}_{n}^{1}a^{n-1}b+\mathrm{C}_{n}^{2}a^{n-2}b^{2}+\mathrm{C}_{n}^{3}a^{n-3}b^{3} + \cdots +  \mathrm{C}_{n}^{n}b^n$.\\
			Cho $a=3$, $b=-1$,ta có 
			\[2^{n}=3^n\mathrm{C}_{n}^{0}-3^{n-1}\mathrm{C}_{n}^{1}+3^{n-2}\mathrm{C}_{n}^{2}-3^{n-3}\mathrm{C}_{n}^{3} +\cdots  +(-1)^n\mathrm{C}_{n}^{n}=2048.\]
			Suy ra $n=10$ vì $n$ chẵn nên số hạng cuối $(-1)^n\mathrm{C}_{n}^{n}=+\mathrm{C}_{10}^{10}$, không  phải $-\mathrm{C}_{10}^{10}$.
			\itemch Ta  có
			$(a+b)^{n}=\mathrm{C}_{n}^{0}a^{n}+\mathrm{C}_{n}^{1}a^{n-1}b+\mathrm{C}_{n}^{2}a^{n-2}b^{2}+\mathrm{C}_{n}^{3}a^{n-3}b^{3} +  \cdots + \mathrm{C}_{n}^{n}b^n$.\\
			Cho $a=1$, $b=1$, $n=12$ ta có 
			 \[2^{n}=\mathrm{C}_{n}^{0}+\mathrm{C}_{n}^{1}+\mathrm{C}_{n}^{2}+\cdots + \mathrm{C}_{n}^{n-1} + \mathrm{C}_{n}^{n}=4096.\]
			Suy ra $\mathrm{C}_{12}^{1}+\mathrm{C}_{12}^{2}  +\cdots+\mathrm{C}_{12}^{11} +\mathrm{C}_{12}^{12}=4095$.
		\end{itemchoice}
	 }
\end{ex}

% nguồn: MaHoaNH_DungSai_Lop10
\begin{ex}%[0D8V3-5]%[ma:c057aeb6]
	Chọn đẳng thức  đúng đúng trong các đẳng thức sau
	 \choiceTF
		 {$\mathrm{C}_{n}^{1}+\mathrm{C}_{n}^{2} +\cdots+\mathrm{C}_{n}^{n-1} +\mathrm{C}_{n}^{n}=2^n$}
		{\True $\mathrm{C}_{2n}^0 + \mathrm{C}_{2n}^2 + \cdots + \mathrm{C}_{2n}^{2n}=\mathrm{C}_{2n}^1 + \mathrm{C}_{2n}^3 +  \cdots + \mathrm{C}_{2n}^{2n-1}=2^{2n-1}$}
		{\True $3^{16}\mathrm{C}_{16}^{0}-3^{15}\mathrm{C}_{16}^{1}+3^{14}\mathrm{C}_{16}^{2}+  \cdots-3\mathrm{C}_{16}^{15}+\mathrm{C}_{16}^{16}=2^{16}$}
		{$\mathrm{C}_{2024}^0 + 3^2\mathrm{C}_{2024}^2 + 3^4\mathrm{C}_{2024}^4 +  \cdots + 3^{2024}\mathrm{C}_{2024}^{2024} = \dfrac{4^{2024}-2^{2024}}{2}$}
	 \loigiai{
		\begin{itemchoice}
			\itemch Ta  có
			$(a+b)^{n}=\mathrm{C}_{n}^{0}a^{n}+\mathrm{C}_{n}^{1}a^{n-1}b+\mathrm{C}_{n}^{2}a^{n-2}b^{2}+\mathrm{C}_{n}^{3}a^{n-3}b^{3} +  \cdots + \mathrm{C}_{n}^{n}b^n$.\\
			Cho  $a=1$, $b=1$ ta có
			\[2^{n}=\mathrm{C}_{n}^{0}+\mathrm{C}_{n}^{1}+\mathrm{C}_{n}^{2}+\cdots + \mathrm{C}_{n}^{n-1}  + \mathrm{C}_{n}^{n}.\]
			Suy ra $\mathrm{C}_{n}^{1}+\mathrm{C}_{n}^{2} +\cdots+\mathrm{C}_{n}^{n-1} +\mathrm{C}_{n}^{n}=2^n-1$. 
			\itemch Xét $ (1+1)^{2n}=\displaystyle\sum\limits_{k=0}^{2n} \mathrm{C}_{2n}^k \cdot  1^{2n-k} \cdot 1^k =\mathrm{C}_{2n}^0 + \mathrm{C}_{2n}^1 + \mathrm{C}_{2n}^2 + \mathrm{C}_{2n}^3 + \mathrm{C}_{2n}^4 + \cdots + \mathrm{C}_{2n}^{2n} $.\tagEX{1} \noindent
			Xét $ (1-1)^{2n}=\displaystyle\sum\limits_{k=0}^{2n} \mathrm{C}_{2n}^k \cdot  1^{2n-k} \cdot (-1)^k =\mathrm{C}_{2n}^0 - \mathrm{C}_{2n}^1 + \mathrm{C}_{2n}^2 - \mathrm{C}_{2n}^3 + \mathrm{C}_{2n}^4 + \cdots + \mathrm{C}_{2n}^{2n} $.\tagEX{2} \noindent
			Lấy $ (1) $ cộng $ (2) $  ta được
			\allowdisplaybreaks 
			\begin{eqnarray*}
				&& 2^{2n} + 0^{2n} =2 \left( \mathrm{C}_{2n}^0 + \mathrm{C}_{2n}^2  + \mathrm{C}_{2n}^4 + \mathrm{C}_{2n}^6 + \cdots + \mathrm{C}_{2n}^{2n} \right) \\
				&\Leftrightarrow& \mathrm{C}_{2n}^0 + \mathrm{C}_{2n}^2 +  \mathrm{C}_{2n}^4 + \mathrm{C}_{2n}^6 + \cdots + \mathrm{C}_{2n}^{2n} =2^{2n-1}.
			\end{eqnarray*}
			Lấy $ (1) $ trừ $ (2) $ ta  được
			\allowdisplaybreaks 
			\begin{eqnarray*}
				&& 2^{2n} - 0^{2n} =2  \left( \mathrm{C}_{2n}^1 + \mathrm{C}_{2n}^3 + \mathrm{C}_{2n}^5 + \mathrm{C}_{2n}^7 + \cdots + \mathrm{C}_{2n}^{2n-1} \right) \\
				&\Leftrightarrow& \mathrm{C}_{2n}^1 + \mathrm{C}_{2n}^3 + \mathrm{C}_{2n}^5  + \mathrm{C}_{2n}^7 + \cdots + \mathrm{C}_{2n}^{2n-1} =2^{2n-1}.
			\end{eqnarray*}
			Vậy $\mathrm{C}_{2n}^0 + \mathrm{C}_{2n}^2  + \cdots + \mathrm{C}_{2n}^{2n}=\mathrm{C}_{2n}^1 + \mathrm{C}_{2n}^3 + \cdots + \mathrm{C}_{2n}^{2n-1}=2^{2n-1}$.
			 \itemch Ta có
			$(a-b)^{16}=\mathrm{C}_{16}^{0}a^{16}-\mathrm{C}_{16}^{1}a^{15}b+\mathrm{C}_{16}^{2}a^{14}b^{2}-\mathrm{C}_{16}^{3}a^{12}b^3+ \cdots  +\mathrm{C}_{16}^{16}b^{16}$.\\
			Cho $a=3$, $b=1$,  ta có
			\[2^{16}=3^{16}\mathrm{C}_{16}^{0}-3^{15}\mathrm{C}_{16}^{1}+3^{14}\mathrm{C}_{16}^{2}+  \cdots-3\mathrm{C}_{16}^{15}+\mathrm{C}_{16}^{16}\]
			 \itemch Xét
			 \allowdisplaybreaks
			\begin{align*}
				 (3+1)^{2024}=&\displaystyle\sum\limits_{k=0}^{2024} \mathrm{C}_{2024}^k \cdot 1^{2024-k} \cdot 3^{k} \\
				 =&\mathrm{C}_{2024}^0 + 3\mathrm{C}_{2024}^1 + 3^2\mathrm{C}_{2024}^2 + 3^3\mathrm{C}_{2024}^3 +3^4 \mathrm{C}_{2024}^4 + \cdots +3^{2024} \mathrm{C}_{2024}^{2024}.\tag{1}
			\end{align*}
			Xét 
			 \allowdisplaybreaks
			\begin{align*}
				(3-1)^{2024}=&\displaystyle\sum\limits_{k=0}^{2024} \mathrm{C}_{2024}^k  \cdot 1^{2024-k} \cdot (-3)^{k} \\
				=&\mathrm{C}_{2024}^0 - 3\mathrm{C}_{2024}^1 + 3^2\mathrm{C}_{2024}^2 - 3^3\mathrm{C}_{2024}^3 +3^4 \mathrm{C}_{2024}^4 - \cdots +3^{2024}  \mathrm{C}_{2024}^{2024}.\tag{2}
			\end{align*}
			Lấy $ (1) $ cộng  $ (2) $ ta được
			\allowdisplaybreaks 
			\begin{eqnarray*}
				&& 4^{2024} + (-2)^{2024} =2 \left(\mathrm{C}_{2024}^0 + 3^2\mathrm{C}_{2024}^2 + 3^4\mathrm{C}_{2024}^4 +  \cdots + 3^{2023}\mathrm{C}_{2024}^{2024} \right) \\
				&\Leftrightarrow&\mathrm{C}_{2024}^0 + 3^2\mathrm{C}_{2024}^2 + 3^4\mathrm{C}_{2024}^4 +  \cdots + 3^{2024}\mathrm{C}_{2024}^{2023} = \dfrac{4^{2024}+2^{2024}}{2}.
			\end{eqnarray*}
		\end{itemchoice}
	} 
\end{ex}


%%%%%%%%%%%%%%%% TRẢ LỜI NGẮN %%%%%%%%%%%%%%%%

% nguồn: MaHoaNH_TraLoiNgan_Lop10
\begin{bt}%[0D8H3-5]%[ma:a490a922]
	Tính tổng $S = \mathrm{C}^{0}_{4}-3\mathrm{C}^{1}_{4} + 9\mathrm{C}^{2}_{4}  -27\mathrm{C}^{3}_{4}+81\mathrm{C}^{4}_{4}$.
	 \shortans{16}
	\loigiai{ 
	Ta  có $S = \mathrm{C}^{0}_{4}(-1)^4 + \mathrm{C}^{1}_{4}(-1)^3\cdot 3 + \mathrm{C}^{2}_{4}(-1)^2\cdot 3^2 + \mathrm{C}^{3}_{4}(-1)\cdot 3^3 + \mathrm{C}^{4}_{4}3^4 = (-1+3)^4 = 16$.
	} 
\end{bt}

% nguồn: MaHoaNH_TraLoiNgan_Lop10
\begin{bt}%[0D8H3-5]%[ma:934d661f]
	 Khai triển $\left(1-\dfrac{1}{2}x\right)^5=a_0+a_1x+a_2x^2+a_3x^3+ a_4x^4+ a_5x^5$. Tính tổng $S=a_0+2a_1+2^2a_2+2^3a_3+2^4a_4+2^5a_5$.
	 \shortans{0}
	 \loigiai{
		Với $x=2$,  ta được
		$$0=\left(1-\dfrac{1}{2}\cdot 2\right)^5=a_0+2a_1+2^2a_2+2^3a_3+2^4a_4+2^5a_5.$$ 
	} 
\end{bt}

% nguồn: MaHoaNH_TraLoiNgan_Lop10
\begin{bt}%[0D8V3-5]%[ma:e6bd0870]
	Lớp $10A$ đề nghị các tổ chọn thành viên để tập kịch. Tổ I phải chọn ít nhất một thành viên để tham gia đội kịch của lớp. Hỏi tổ I có bao nhiêu cách chọn thành viên để tập kịch? Biết rằng tổ I có $5$ người. 
	\shortans{31} 
	 \loigiai{
		Vì tổ I phải chọn ít nhất  một thành viên để tham gia đội kịch nên số cách chọn thành viên của tổ I là
		\begin{eqnarray*}
			&&\mathrm{C}_5^1+\mathrm{C}_5^2+\mathrm{C}_5^3+\mathrm{C}_5^4+\mathrm{C}_5^5=(1+1)^5-\mathrm{C}_5^0=2^5-1=31. 
		\end{eqnarray*}
	} 
\end{bt}

% nguồn: MaHoaNH_TraLoiNgan_Lop10
\begin{bt}%[0D8V3-5]%[ma:7d215490]
	Tìm số nguyên dương $m$  biết $\mathrm{C}^{0}_{20} + \mathrm{C}^{0}_{20} + \cdots + \mathrm{C}^{20}_{20} = 2^m$.
	\shortans{19} 
	 \loigiai{
	Xét khai triển $f(x) = (1+x)^{20} = \mathrm{C}^{0}_{20} + \mathrm{C}^{1}_{20}x + \mathrm{C}^{2}_{20}x^2 + \cdots  + \mathrm{C}^{20}_{20}x^{20}$.\\
	Do $f(-1) =  0$ nên
	 $$
	\mathrm{C}^{0}_{20} - \mathrm{C}^{1}_{20} + \mathrm{C}^{2}_{20} -  \cdots + \mathrm{C}^{20}_{20} = 0
	\Leftrightarrow \mathrm{C}^{0}_{20} + \mathrm{C}^{0}_{20} + \cdots + \mathrm{C}^{20}_{20} = \mathrm{C}^{1}_{20} +  \mathrm{C}^{3}_{20} + \cdots + \mathrm{C}^{19}_{20}.
	 $$
	Kết  hợp với $f(1) = 2^{20}$ ta được
	$$ 
	2^{20} = \mathrm{C}^{0}_{20} + \mathrm{C}^{1}_{20} + \cdots + \mathrm{C}^{20}_{20} = 2\left(\mathrm{C}^{0}_{20} + \mathrm{C}^{0}_{20} + \cdots +  \mathrm{C}^{20}_{20}\right) = 2^{m+1} \Leftrightarrow m=19.
	$$ 
	Vậy  $m=19$.
	 }
\end{bt}

% nguồn: MaHoaNH_TraLoiNgan_Lop10
\begin{bt}%[0D8C3-5]%[ma:c52f2c6b]
	Cho $n$ là số tự nhiên thỏa mãn $\mathrm{C}^{0}_{n}+2\mathrm{C}^{1}_{n}+2^2\mathrm{C}^{2}_{n}+\cdots+2^n\mathrm{C}^{n}_{n} = 59049$. Biết số hạng thứ $3$ trong khai triển Newton của $\left(x^2-\dfrac{3}{x}\right)^n$ có giá trị bằng $\dfrac{81n}{2}$, tìm số  giá trị của $x$ thỏa mãn.
	\shortans{2} 
	\loigiai{ 
	Ta có $\mathrm{C}^{0}_{n}+2\mathrm{C}^{1}_{n}+2^2\mathrm{C}^{2}_{n}+\cdots+2^n\mathrm{C}^{n}_{n} = (1+2)^n =  3^n$ nên $3^n = 59049$ hay $n=10$.\\
	Từ đó ta được nhị thức $\left(x^2 -  \dfrac{3}{x}\right)^{10}$ có số hạng tổng quát là
	$$ 
	T_{k+1} = \mathrm{C}^{k}_{10}(x^2)^{10-k}\left(-\dfrac{3}{x}\right)^k = \mathrm{C}^{k}_{10}(-3)^kx^{20-3k} \text{ với } k  = \overline{0,1,\ldots,10}.
	 $$
	Khi đó số hạng thứ $3$  là $T_3 = \mathrm{C}^{2}_{10}(-3)^2x^{14} = 405x^{14}$.\\
	Theo bài ra ta có $405x^{14} = \dfrac{81n}{2} = 405 \Leftrightarrow x^{14}  = 1 \Leftrightarrow x=\pm 1$.\\
	Vậy có $2$ giá trị $x$ thỏa mãn. 
	 }
\end{bt}


%%%%%%%%%%%%%%%% TỰ LUẬN %%%%%%%%%%%%%%%%

% nguồn: MaHoaNH_TuLuan_Lop10
\begin{bt}%[0D8H3-5]%[ma:18c0055d]
	Cho  đa thức $P(x)=(1-x)^8$. Tính tổng các hệ số của đa thức $P(x)$.
	\loigiai{ 
		Ta có  $P(x)=(1-x)^8=\displaystyle\displaystyle\sum\limits_{k=0}^8 \mathrm{C}_8^k(-1)^k x^k$.\\
		Khi đó tổng các hệ số của đa thức  $P(x)$ là
		 $S=\mathrm{C}_8^0-\mathrm{C}_8^1+\cdots-\mathrm{C}_8^7+\mathrm{C}_8^8$.\\
		Xét khai triển $(a+b)^8=\displaystyle\sum\limits_{k=0}^8  \mathrm{C}_8^k a^{8-k} b^k$.\\
		Ta chọn  $a=1$, $b=-1$ thu được $(1-1)^8=\mathrm{C}_8^0-\mathrm{C}_8^1+\mathrm{C}_8^2-\cdots-\mathrm{C}_8^7+\mathrm{C}_8^8$.\\
		Vậy  tổng các hệ số của đa thức $P(x)$ bằng $0$.
	 }
\end{bt}

% nguồn: MaHoaNH_TuLuan_Lop10
\begin{bt}%[0D8H3-5]%[ma:6d7a94f0]
	Tính tổng $S=2^7 \mathrm{C}_{2n}^0-2^8 \mathrm{C}_{2n}^1+2^9  \mathrm{C}_{2n}^2-2^{10} \mathrm{C}_{2n}^3+\cdots-2^{2n+6} \mathrm{C}_{2n}^{2n-1}+2^{2n+7} \mathrm{C}_{2n}^{2n}$ với $n \in \mathbb{N}^*$.
	\loigiai{ 
		Ta có $S=2^7\left(\mathrm{C}_{2n}^0-2^1 \mathrm{C}_{2n}^1+2^2 \mathrm{C}_{2n}^2-2^3  \mathrm{C}_{2n}^3+\cdots-2^{2n-1} \mathrm{C}_{2n}^{2n-1}+2^{2n} \mathrm{C}_{2n}^{2n}\right)$.\\
		Xét khai triển 
		\begin{eqnarray*}
			(x-2)^{2n}=\mathrm{C}_{2n}^0 x^{2n}(-2)^0 +\mathrm{C}_{2n}^1  x^{2n-1} \cdot(-2)^1 &+&\mathrm{C}_{2n}^2 x^{2n-2} \cdot(-2)^2+\cdots\\
			 &+&\mathrm{C}_{2n}^{2n-1} x^1(-2)^{2n-1}+\mathrm{C}_{2n}^{2n}(-2)^{2n}.
		\end{eqnarray*}
		Chọn $x=1$ ta có $1=(-1)^{2n}=\mathrm{C}_{2n}^0-2^1 \mathrm{C}_{2n}^1+2^2 \mathrm{C}_{2n}^2-2^3 \mathrm{C}_{2n}^3+\cdots-2^{2n-1}  \mathrm{C}_{2n}^{2n-1}+2^{2n} \mathrm{C}_{2n}^{2n}$.\\
		 Vậy $S=2^7 (-1)^{2n}=2^7$.
	} 
\end{bt}

% nguồn: MaHoaNH_TuLuan_Lop10
\begin{bt}%[0D8V3-5]%[ma:ba23ea5e]
	Rút gọn biểu thức $S=\dfrac{1}{1\cdot 0!\cdot 2019!}+\dfrac{1}{2\cdot 1!\cdot  2018!}+\dfrac{1}{3\cdot 2!\cdot 2017!}+\cdots+\dfrac{1}{2020\cdot 2019!\cdot 0!}$.
	 \loigiai{
		Ta có  $S=\displaystyle\sum\limits_{k=0}^{2019} \dfrac{1}{(k+1) k!(2019-k)!}=\displaystyle\sum\limits_{k=0}^{2019} \dfrac{2020!}{2020!(k+1)!(2020-(k+1))!}=\dfrac{1}{2020!} \displaystyle\sum\limits_{k=0}^{2019} \mathrm{C}_{2020}^{k+1}$.\\
		Xét khai triển $(x+1)^{2020}=\displaystyle\sum\limits_{k=0}^{2020} \mathrm{C}_{2020}^k \cdot x^k=1+\displaystyle\sum\limits_{k=1}^{2020} \mathrm{C}_{2020}^k \cdot x^k$.\\ 
		Chọn $x=1$ ta thu được $\displaystyle\sum\limits_{k=1}^{2020} \mathrm{C}_{2020}^k=\displaystyle\sum\limits_{k=0}^{2019}  \mathrm{C}_{2020}^{k+1}=2^{2020}-1$.\\
		 Vậy $S=\dfrac{2^{2020}-1}{2020!}$.
	 }
\end{bt}

% nguồn: MaHoaNH_TuLuan_Lop10
\begin{bt}%[0D8V3-5]%[ma:9a6ef033]
	Cho $n$ là số tự nhiên. Thu gọn biểu thức $S=3 \mathrm{C}_n^0+7 \mathrm{C}_n^1+11 \mathrm{C}_n^2+\cdots+(4n+3) \mathrm{C}_n^n$ theo  $n$.
	 \loigiai{
		Ta  có $S=(0\cdot 4+3) \mathrm{C}_n^0+(1\cdot 4+3) \mathrm{C}_n^1+(2\cdot 4+3) \mathrm{C}_n^2+\cdots+(n\cdot 4+3) \mathrm{C}_n^n$.\\
		Suy ra $S=4\left(\mathrm{C}_n^1+2 \mathrm{C}_n^2+3 \mathrm{C}_n^3+\cdots+n  \cdot \mathrm{C}_n^n\right)+3\left(\mathrm{C}_n^0+\mathrm{C}_n^1+\cdots+\mathrm{C}_n^n\right)$.\\
		Xét khai triển $(x+1)^n=\mathrm{C}_n^0 x^0+\mathrm{C}_n^1 x^1+\cdots+\mathrm{C}_n^n  x^n$.\\
		Chọn $x=1$ ta thu được  $\mathrm{C}_n^0+\mathrm{C}_n^1+\cdots+\mathrm{C}_n^n=2^n$.\\
		Mặt khác ta lại có $k \cdot \mathrm{C}_n^k=k \cdot \dfrac{n!}{k!(n-k)!}=\dfrac{n(n-1)!}{(k-1)![(n-1)-(k-1)]!}=n  \cdot \mathrm{C}_{n-1}^{k-1}$.\\
		Do đó $\mathrm{C}_n^1+2 \cdot \mathrm{C}_n^2+3\cdot  \mathrm{C}_n^3+\cdots+n \cdot \mathrm{C}_n^n=n\left(\mathrm{C}_{n-1}^0+\mathrm{C}_{n-1}^1+\mathrm{C}_{n-1}^2+\cdots+\mathrm{C}_{n-1}^{n-1}\right)$.\\
		Tương  tự xét khai triển $(x+1)^{n-1}=\mathrm{C}_{n-1}^0 \cdot x^0+\mathrm{C}_{n-1}^1 \cdot x^1+\cdots+\mathrm{C}_{n-1}^{n-1} \cdot x^{n-1}$.\\
		 Chọn $x=1$ ta thu được $\mathrm{C}_{n-1}^0+\mathrm{C}_{n-1}^1+\mathrm{C}_{n-1}^2+\cdots+\mathrm{C}_{n-1}^{n-1}=2^{n-1}$.\\
		Vậy $S=4n \cdot 2^{n-1}+3 \cdot  2^n=(2n+3) \cdot 2^n$.
	 }
\end{bt}

% nguồn: MaHoaNH_TuLuan_Lop10
\begin{bt}%[0D8V3-5]%[ma:badfa851]
	Tính  tổng của các đẳng thức sau
	\begin{listEX}
		 \item $S=\mathrm{C}_{10}^0+\mathrm{C}_{10}^1+\cdots+\mathrm{C}_{10}^{10}$.
		\item $S=\mathrm{C}_6^1+\mathrm{C}_6^2+\cdots+\mathrm{C}_6^5$. 
		\item $S=\mathrm{C}_6^0+2  \cdot \mathrm{C}_6^1+2^2 \cdot \mathrm{C}_6^2+\cdots+2^6 \mathrm{C}_6^6$.
		\item $S=\mathrm{C}_{12}^0-\mathrm{C}_{12}^1+\mathrm{C}_{12}^2-\cdots-\mathrm{C}_{12}^{11}+\mathrm{C}_{12}^{12}$. 
		\item $S=\mathrm{C}_{20}^1+2 \mathrm{C}_{20}^2+2^2  \cdot \mathrm{C}_{20}^3+\cdots+2^{19} \mathrm{C}_{20}^{20}$.
		 \item $S=\mathrm{C}_{20}^0+\mathrm{C}_{20}^2+\mathrm{C}_{20}^4+\cdots+\mathrm{C}_{20}^{20}$.
		\item $S=\mathrm{C}_{2019}^1 \cdot 3^{2018} \cdot 2-\mathrm{C}_{2019}^2 \cdot 3^{2017} \cdot 2^2+\mathrm{C}_{2019}^3  \cdot 3^{2016} \cdot 2^3-\cdots-\mathrm{C}_{2019}^{2018} \cdot 3^1 \cdot 2^{2018}+\mathrm{C}_{2019}^{2019} \cdot 2^{2019}$.
		\item $S=\mathrm{C}_{2021}^0 \cdot 4^{2021}-\mathrm{C}_{2021}^1 \cdot 4^{2010} \cdot 2+\mathrm{C}_{2021}^2 \cdot 4^{2019} \cdot  2^2-\mathrm{C}_{2021}^3 \cdot 4^{2018} \cdot 2^3-\cdots+\mathrm{C}_{2021}^{2020} \cdot 4^1 \cdot 2^{2020}$.
	\end{listEX}
	\loigiai{ 
		\begin{listEX}
			\item $S=\mathrm{C}_{10}^0+\mathrm{C}_{10}^1+\cdots+\mathrm{C}_{10}^{10}$.\\ 
			Xét  khai triển $(a+b)^{10}=\displaystyle\sum\limits_{k=0}^{10} \mathrm{C}_{10}^k a^{10-k} b^k$.\\
			Ta chọn $a=b=1$, thu được  $(1+1)^{10}=\mathrm{C}_{10}^0+\mathrm{C}_{10}^1+\cdots+\mathrm{C}_{10}^{10}$ nên $S=2^{10}=1\,024$.
			 \item $S=\mathrm{C}_6^1+\mathrm{C}_6^2+\cdots+\mathrm{C}_6^5$.\\
			Xét khai triển $(a+b)^6=\displaystyle\sum\limits_{k=0}^6 \mathrm{C}_6^k a^{6-k}  b^k$.\\
			Ta chọn $a=b=1$, thu  được $(1+1)^6=\mathrm{C}_6^0+\mathrm{C}_6^1+\cdots+\mathrm{C}_6^6$ nên $S=2^6-\mathrm{C}_6^0-\mathrm{C}_6^6=62$.
			\item $S=\mathrm{C}_6^0+2 \cdot  \mathrm{C}_6^1+2^2 \cdot \mathrm{C}_6^2+\cdots+2^6 \mathrm{C}_6^6$.\\
			Xét khai triển $(a+b)^6=\displaystyle\sum\limits_{k=0}^6 \mathrm{C}_6^k a^{6-k}  b^k$.\\
			Ta chọn $a=1$, $b=2$ thu được $(1+2)^6=\mathrm{C}_6^0+2  \cdot \mathrm{C}_6^1+2^2 \cdot \mathrm{C}_6^2+\cdots+2^6 \mathrm{C}_6^6$ nên $S=3^6=729$.
			 \item $S=\mathrm{C}_{12}^0-\mathrm{C}_{12}^1+\mathrm{C}_{12}^2-\cdots-\mathrm{C}_{12}^{11}+\mathrm{C}_{12}^{12}$.\\
			 Xét khai triển $(a+b)^{12}=\displaystyle\sum\limits_{k=0}^{12} \mathrm{C}_{12}^k a^{12-k} b^k$.\\
			Ta chọn $a=1$, $b=-1$ thu được $(1-1)^{12}=\mathrm{C}_{12}^0-\mathrm{C}_{12}^1+\mathrm{C}_{12}^2-\cdots-\mathrm{C}_{12}^{11}+\mathrm{C}_{12}^{12}$ nên  $S=0^{12}=0$.
			\item $S=\mathrm{C}_{20}^1+2 \mathrm{C}_{20}^2+2^2  \cdot \mathrm{C}_{20}^3+\cdots+2^{19} \mathrm{C}_{20}^{20}$.\\
			Ta có $2 S=2 \cdot \mathrm{C}_{20}^1+2^2 \mathrm{C}_{20}^2+2^3 \cdot  \mathrm{C}_{20}^3+\cdots+2^{20} \cdot \mathrm{C}_{20}^{20}$.\\
			Xét khai triển $(a+b)^{20}=\displaystyle\sum\limits_{k=0}^{20} \mathrm{C}_{20}^k a^{20-k} b^k$.\\ 
			Ta chọn $a=1$, $b=2$ thu được  $(1+2)^{20}=\mathrm{C}_{20}^0+2 \cdot \mathrm{C}_{20}^1+\cdots+2^{20} \cdot \mathrm{C}_{20}^{20}$.\\
			Do đó  $2 S=(1+2)^{20}-\mathrm{C}_{20}^0=3^{20}-1$ nên $S=\dfrac{3^{20}-1}{2}$.
			\item  $S=\mathrm{C}_{20}^0+\mathrm{C}_{20}^2+\mathrm{C}_{20}^4+\cdots+\mathrm{C}_{20}^{20}$.\\
			Xét khai triển $(a+b)^{20}=\displaystyle\sum\limits_{k=0}^{20}  \mathrm{C}_{20}^k a^{20-k} b^k$.\\
			 Chọn $a=b=1$, ta thu được $(1+1)^{20}=\mathrm{C}_{20}^0+\mathrm{C}_{20}^1+\mathrm{C}_{20}^2+\mathrm{C}_{20}^3 +\cdots+\mathrm{C}_{20}^{20}$.\\
			Chọn  $a=1$, $b=-1$ ta thu được $(1-1)^{20}=\mathrm{C}_{20}^0-\mathrm{C}_{20}^1+\mathrm{C}_{20}^2-\mathrm{C}_{20}^3+\cdots+\mathrm{C}_{20}^{20}$.\\
			Cộng theo vế  hai phương trình ta được
			$$2^{20}=2 \cdot\left(\mathrm{C}_{20}^0+\mathrm{C}_{20}^2+\mathrm{C}_{20}^4+\cdots+\mathrm{C}_{20}^{20}\right) \Leftrightarrow 2 S=2^{20}  \Leftrightarrow S=2^{19}.$$
			\item $S=\mathrm{C}_{2019}^1 \cdot 3^{2018} \cdot  2-\mathrm{C}_{2019}^2 \cdot 3^{2017} \cdot 2^2+\mathrm{C}_{2019}^3 \cdot 3^{2016} \cdot 2^3-\cdots-\mathrm{C}_{2019}^{2018} \cdot 3^1 \cdot 2^{2018}+\mathrm{C}_{2019}^{2019} \cdot 2^{2019}$.\\
			Xét khai triển \begin{eqnarray*}
			&&(a+b)^{2019} 
			=\displaystyle\sum\limits_{k=0}^{2019} \mathrm{C}_{2019}^k a^{2019-k} b^k\\ 
			&=&\mathrm{C}_{2019}^0 \cdot a^{2019}+\mathrm{C}_{2019}^1 \cdot a^{2018} \cdot b+\mathrm{C}_{2019}^2 \cdot a^{2017} \cdot b^2+\mathrm{C}_{2019}^3 \cdot  a^{2016} \cdot b^3+\cdots\\
			&&\quad\quad\quad\quad\quad\quad\quad\quad\quad\quad\quad\quad\quad\quad\quad\quad\quad+\mathrm{C}_{2019}^{2018} \cdot a^1 \cdot  b^{2018}+\mathrm{C}_{2019}^{2019} \cdot b^{2019}.
		\end{eqnarray*}
		Ta chọn  $a=-3$, $b=2$ khi đó
		\begin{eqnarray*}
			(-3+2)^{2019}=-\mathrm{C}_{2019}^0 \cdot 3^{2019}+\mathrm{C}_{2019}^1 \cdot 3^{2018} \cdot  2 &-&\mathrm{C}_{2019}^2 \cdot 3^{2017} \cdot 2^2
			+\mathrm{C}_{2019}^3  \cdot 3^{2016} \cdot 2^3+\cdots\\
			&-&\mathrm{C}_{2019}^{2018} \cdot 3^1 \cdot 2^{2018}+\mathrm{C}_{2019}^{2019}  \cdot 2^{2019}.
		\end{eqnarray*}
		Mà $S=\mathrm{C}_{2019}^1 \cdot 3^{2018} \cdot 2-\mathrm{C}_{2019}^2 \cdot  3^{2017} \cdot 2^2+\mathrm{C}_{2019}^3 \cdot 3^{2016} \cdot 2^3+\cdots-\mathrm{C}_{2019}^{2018} \cdot 3^1 \cdot 2^{2018}+\mathrm{C}_{2019}^{2019} \cdot 2^{2019}$.\\
		 Vậy $S=(-3+2)^{2019}+\mathrm{C}_{2019}^0 \cdot 3^{2019}=(-1)^{2019}+3^{2019}=3^{2019}-1$.
		\item $S=\mathrm{C}_{2021}^0 \cdot 4^{2021}-\mathrm{C}_{2021}^1 \cdot 4^{2010} \cdot 2+\mathrm{C}_{2021}^2 \cdot 4^{2019} \cdot 2^2-\mathrm{C}_{2021}^3 \cdot 4^{2018} \cdot 2^3-\cdots+\mathrm{C}_{2021}^{2020} \cdot 4^1 \cdot  2^{2020}$.
		Xét khai  triển
		\begin{eqnarray*}
			&&(a+b)^{2021}=\displaystyle\sum\limits_{k=0}^{2021} \mathrm{C}_{2021}^k  1^{2021-k} b^k\\
			&=&\mathrm{C}_{2021}^0 \cdot a^{2021}+\mathrm{C}_{2021}^1 \cdot a^{2020} \cdot b+\mathrm{C}_{2021}^2 \cdot a^{2019} \cdot  b^2+\mathrm{C}_{2021}^3 \cdot a^{2018} \cdot b^3+\cdots\\
			&&\quad\quad\quad\quad\quad\quad\quad\quad\quad\quad\quad\quad\quad\quad\quad\quad\quad+\mathrm{C}_{2021}^{2020} \cdot a^1  \cdot b^{2020}+\mathrm{C}_{2021}^{2021} \cdot b^{2021}.
		\end{eqnarray*}
		Ta chọn $a=4$,  $b=-2$ khi đó
		\begin{eqnarray*}
			&&(4-2)^{2021}\\ 
			&=&\mathrm{C}_{2021}^0 \cdot 4^{2021}-\mathrm{C}_{2021}^1 \cdot 4^{2020} \cdot 2+\mathrm{C}_{2021}^2 \cdot 4^{2019} \cdot 2^2-\mathrm{C}_{2021}^3  \cdot 4^{2018} \cdot 2^3+\cdots\\
			&&\quad\quad\quad\quad\quad\quad\quad\quad\quad\quad\quad\quad\quad\quad\quad\quad\quad\quad+\mathrm{C}_{2021}^{2020} \cdot 4  \cdot 2^{2020}
			-\mathrm{C}_{2021}^{2021} \cdot  2^{2021}.
		\end{eqnarray*}
		Mà $S=\mathrm{C}_{2021}^0 \cdot 4^{2021}-\mathrm{C}_{2021}^1 \cdot 4^{2020} \cdot 2+\mathrm{C}_{2021}^2 \cdot 4^{2019} \cdot 2^2-\mathrm{C}_{2021}^3 \cdot 4^{2018}  \cdot 2^3+\cdots+\mathrm{C}_{2021}^{2020} \cdot 4 \cdot 2^{2020}$.\\
		Vậy  $S=(4-2)^{2021}+\mathrm{C}_{2021}^{2021} \cdot 2^{2021}=2^{2021}+2^{2021}=2^{2022}$.
	\end{listEX}
 }
\end{bt}
